%% lecture 2 %%%%
\documentclass[dvipdfmx,a4paper]{jsarticle}
\usepackage{amsmath,amssymb,amsfonts,amsthm}
\usepackage[dvipdfmx]{graphicx}
\usepackage{tikz}
\usetikzlibrary{positioning, intersections, calc, arrows.meta, math, through, shadows,automata}
\usepackage{caption}
\usepackage{tcolorbox}
\tcbuselibrary{theorems,breakable}
\usepackage{enumerate}
\usepackage{mathtools}
\usepackage{otf}
\usepackage{xspace}
\usepackage{newpxtext}
\usepackage[utf8]{inputenc}
\usepackage{CJKutf8,CJKspace,CJKpunct}
\usepackage{pgfplots}
\usepackage{lastpage}
\usepackage{fancyhdr}
\usepackage{float}
\usepackage{lipsum}
\usepackage{okumacro}
\usepackage[dvipdfmx,colorlinks=true, linkcolor=blue]{hyperref}

\pagestyle{fancy}
\fancyhf{}
\fancyhead[L]{\textsc{特殊講義（専門01)}}
\fancyhead[R]{\texttt{YI Ran} - $\mathnormal{61112600301}$}
\fancyfoot[C]{\thepage\quad \textit{of}\quad\pageref*{LastPage}}


\fancypagestyle{plain}{
\fancyhf{}
\renewcommand{\headrulewidth}{0pt}
\fancyfoot[C]{\thepage\quad \textit{of}\quad\pageref*{LastPage}}
}

\pgfplotsset{compat=1.18}
\renewcommand{\appendixname}{Appendix}

\renewcommand{\abstractname}{問題設定}

\newtagform{textbf}[\textbf]{[}{]}
\usetagform{textbf}

\newcommand*{\ie}{\textbf{\textit{i.e.}}\@\xspace}

\renewcommand{\qedsymbol}{$\blacksquare$}

\newtcbtheorem[]{reidai}{例題}
{fonttitle=\gtfamily\sffamily\bfseries\upshape\large,
colframe=black,colback=black!15!white,
rightrule=1pt,leftrule=1pt,bottomrule=2pt,
colbacktitle=black,theorem style=standard,breakable,arc=10pt}
{tha}

\renewcommand{\thefootnote}{\arabic{footnote}}

\newtheoremstyle{mystyle}%
{}{}{}{}{\bfseries}{:}{ }
{\thmname{#1}\thmnumber{ #2}\thmnote{ (#3)}}

\theoremstyle{mystyle}

\newtheorem{dfn}{\texttt{Def.}}[section]
\newtheorem{exm}[dfn]{\texttt{Ex.}}
\newtheorem{prop}[dfn]{\texttt{Prop.}}
\newtheorem{lem}[dfn]{\texttt{Lem.}}
\newtheorem{thm}[dfn]{\texttt{Thm.}}
\newtheorem{cor}[dfn]{\texttt{Cor.}}
\newtheorem{rem}[dfn]{\texttt{Rem.}}
\newtheorem{fact}[dfn]{\texttt{Fact}}

\newcommand{\kai}
{\noindent
\begin{tikzpicture}[scale=0.2, baseline=2.8pt]
\draw (3.3,1.2) node{\large\textgt{解 答}};
\draw[thick, rounded corners=3pt,] (0,0)--(6.5,0)--(6.5,2.4)--(0,2.4)--cycle;
\end{tikzpicture}}

\newcommand{\shomei}
{\noindent
\begin{tikzpicture}[scale=0.2, baseline=2.8pt]
\draw (3.3,1.2) node{\textgt{証 明}};
\draw[double,thick,rounded corners=3pt,] (0,0)--(6.5,0)--(6.5,2.4)--(0,2.4)--cycle;
\end{tikzpicture};}

\newcommand{\hosoku}{\noindent
\begin{tikzpicture}[scale=0.2, baseline=2.8pt]
\draw (6,1) node{\large\textgt{補足}};
\fill (0,1)--(1,0)--(2,1)--(1,2)--cycle;
\fill[gray] (1,1)--(2,0)--(3,1)--(2,2)--cycle;
\fill (2,1)--(3,0)--(4,1)--(3,2)--cycle;
\fill (10,1)--(11,0)--(12,1)--(11,2)--cycle;
\fill[gray] (9,1)--(10,0)--(11,1)--(10,2)--cycle;
\fill (8,1)--(9,0)--(10,1)--(9,2)--cycle;
\end{tikzpicture};}

\newcommand{\chui}{\noindent
\begin{tikzpicture}[scale=0.2, baseline=2.8pt]
\fill (0,0)--(6.5,0)--(6.5,2.2)--(0,2.2);
\draw (3.3,1) node[white]{\large\textgt{注意！}};
\draw[thick] (0,0)--(6.5,0)--(6.5,2.2)--(0,2.2)--cycle;
\end{tikzpicture};}

\title{\vspace{-3cm} \textbf{\texttt{REPORT 2(SPECIAL LECTURES)}}}
\author{\texttt{YI Ran} - $\mathnormal{61112600301}$\\ \texttt{andreyi@outlook.jp}}
\date{\today}

\begin{document}
\maketitle
\thispagestyle{plain}

\begin{abstract} %概要
  \noindent
\textrm{Let $H=\mathbb{C}^2$ and $\sigma_0,\sigma_1,\sigma_2$ and $\sigma_3$ be the Pauli matrices on $H$.
That is,}
\[
\sigma_0=
\begin{pmatrix}
1 & 0\\
0 & 1
\end{pmatrix},
\qquad
\sigma_1=
\begin{pmatrix}
0 & 1\\
1 & 0
\end{pmatrix},
\]
\[
\sigma_2=
\begin{pmatrix}
0 & -i\\
i & 0
\end{pmatrix},
\qquad
\sigma_3=
\begin{pmatrix}
1 & 0\\
0 & -1
\end{pmatrix}.
\]
\end{abstract}

%============================================================
\section*{\boxed{\textbf{\large 問1}}}
%============================================================

\textrm{Let $H=\mathbb{C}^2$ and let
\[
|0\rangle=
\begin{pmatrix}
1\\
0
\end{pmatrix},
\qquad
|1\rangle=
\begin{pmatrix}
0\\
1
\end{pmatrix}
\]
be the canonical basis for $\mathbb{C}^2$.
Express $\sigma_1,\sigma_2,\sigma_3$ by the ket vectors $|0\rangle$ and $|1\rangle$.}

\vspace{1em}

%============================================================
\section*{\boxed{\textbf{\large 解答}}}
%============================================================

\begin{proof}
\ \\
まず，$|0\rangle,|1\rangle$ の bra を書くと
\[
\langle 0|=
\begin{pmatrix}
1 & 0
\end{pmatrix},
\qquad
\langle 1|=
\begin{pmatrix}
0 & 1
\end{pmatrix}
\]
である．

\medskip
\noindent
次に，これらから作られる外積を一つずつ計算する．
\[
|0\rangle\langle 0|
=
\begin{pmatrix}
1\\
0
\end{pmatrix}
\begin{pmatrix}
1 & 0
\end{pmatrix}
=
\begin{pmatrix}
1 & 0\\
0 & 0
\end{pmatrix},
\qquad
|0\rangle\langle 1|
=
\begin{pmatrix}
1\\
0
\end{pmatrix}
\begin{pmatrix}
0 & 1
\end{pmatrix}
=
\begin{pmatrix}
0 & 1\\
0 & 0
\end{pmatrix},
\]
\[
|1\rangle\langle 0|
=
\begin{pmatrix}
0\\
1
\end{pmatrix}
\begin{pmatrix}
1 & 0
\end{pmatrix}
=
\begin{pmatrix}
0 & 0\\
1 & 0
\end{pmatrix},
\qquad
|1\rangle\langle 1|
=
\begin{pmatrix}
0\\
1
\end{pmatrix}
\begin{pmatrix}
0 & 1
\end{pmatrix}
=
\begin{pmatrix}
0 & 0\\
0 & 1
\end{pmatrix}.
\]

\medskip
\noindent
以上を用いて $\sigma_1,\sigma_2,\sigma_3$ を表す．

\paragraph{\textbf{$\sigma_1$ の場合}} \ \\
\[
|0\rangle\langle 1|+|1\rangle\langle 0|
=
\begin{pmatrix}
0 & 1\\
0 & 0
\end{pmatrix}
+
\begin{pmatrix}
0 & 0\\
1 & 0
\end{pmatrix}
=
\begin{pmatrix}
0 & 1\\
1 & 0
\end{pmatrix}
=
\sigma_1.
\]
したがって，
\[
\sigma_1=|0\rangle\langle 1|+|1\rangle\langle 0|
\]

\paragraph{\textbf{$\sigma_2$ の場合}} \ \\
\[
-i|0\rangle\langle 1|+i|1\rangle\langle 0|
=
-i
\begin{pmatrix}
0 & 1\\
0 & 0
\end{pmatrix}
+i
\begin{pmatrix}
0 & 0\\
1 & 0
\end{pmatrix}
=
\begin{pmatrix}
0 & -i\\
i & 0
\end{pmatrix}
=
\sigma_2.
\]
したがって，
\[
\sigma_2=-i|0\rangle\langle 1|+i|1\rangle\langle 0|
\]

\paragraph{\textbf{$\sigma_3$ の場合}} \ \\
\[
|0\rangle\langle 0|-|1\rangle\langle 1|
=
\begin{pmatrix}
1 & 0\\
0 & 0
\end{pmatrix}
-
\begin{pmatrix}
0 & 0\\
0 & 1
\end{pmatrix}
=
\begin{pmatrix}
1 & 0\\
0 & -1
\end{pmatrix}
=
\sigma_3.
\]
したがって，
\[
\sigma_3=|0\rangle\langle 0|-|1\rangle\langle 1|
\]

\medskip
\noindent
以上より，$\sigma_1,\sigma_2,\sigma_3$ はそれぞれ
\[
\sigma_1=|0\rangle\langle 1|+|1\rangle\langle 0|,\qquad
\sigma_2=-i|0\rangle\langle 1|+i|1\rangle\langle 0|,\qquad
\sigma_3=|0\rangle\langle 0|-|1\rangle\langle 1|
\]
と表される．
\end{proof}

\newpage

%============================================================
\section*{\boxed{\textbf{\large 問2}}}
%============================================================

Let $H=\mathbb{C}^2\otimes\mathbb{C}^2$.
Show that
\[
|0\rangle\otimes|0\rangle,\quad
|0\rangle\otimes|1\rangle,\quad
|1\rangle\otimes|0\rangle,\quad
|1\rangle\otimes|1\rangle
\]
are the standard basis for $H$.

\vspace{1em}

%============================================================
\section*{\boxed{\textbf{\large 解答}}}
%============================================================

\begin{proof}
\ \\
$\mathbb{C}^2$ の標準基底を
\[
|0\rangle=
\begin{pmatrix}
1\\
0
\end{pmatrix},
\qquad
|1\rangle=
\begin{pmatrix}
0\\
1
\end{pmatrix}
\]
とする。$\dim{H} = \dim{(\mathbb{C}^2\otimes\mathbb{C}^2)}=\dim{\mathbb{C}^2}\cdot \dim{\mathbb{C}^2}=2\cdot 2=4 = \dim{\mathbb{C}^4}$である。\ie $H\cong\mathbb{C}^4$\\
\noindent
(詳しい証明は \texttt{\ref{proof:dim}} を参照)。

\medskip
\noindent
次、ベクトルのテンソル積を計算する。
\paragraph{\textbf{$|0\rangle\otimes|0\rangle$ の計算}} \ \\
\[
|0\rangle\otimes|0\rangle
=
\begin{pmatrix}
1\\
0
\end{pmatrix}
\otimes
\begin{pmatrix}
1\\
0
\end{pmatrix}
=
\begin{pmatrix}
1\cdot 1\\
1\cdot 0\\
0\cdot 1\\
0\cdot 0
\end{pmatrix}
=
\begin{pmatrix}
1\\
0\\
0\\
0
\end{pmatrix}
\]

\paragraph{\textbf{$|0\rangle\otimes|1\rangle$ の計算}} \ \\
\[
|0\rangle\otimes|1\rangle
=
\begin{pmatrix}
1\\
0
\end{pmatrix}
\otimes
\begin{pmatrix}
0\\
1
\end{pmatrix}
=
\begin{pmatrix}
1\cdot 0\\
1\cdot 1\\
0\cdot 0\\
0\cdot 1
\end{pmatrix}
=
\begin{pmatrix}
0\\
1\\
0\\
0
\end{pmatrix}
\]

\paragraph{\textbf{$|1\rangle\otimes|0\rangle$ の計算}} \ \\
\[
|1\rangle\otimes|0\rangle
=
\begin{pmatrix}
0\\
1
\end{pmatrix}
\otimes
\begin{pmatrix}
1\\
0
\end{pmatrix}
=
\begin{pmatrix}
0\cdot 1\\
0\cdot 0\\
1\cdot 1\\
1\cdot 0
\end{pmatrix}
=
\begin{pmatrix}
0\\
0\\
1\\
0
\end{pmatrix}
\]

\paragraph{\textbf{(iv) $|1\rangle\otimes|1\rangle$ の計算}} \ \\
\[
|1\rangle\otimes|1\rangle
=
\begin{pmatrix}
0\\
1
\end{pmatrix}
\otimes
\begin{pmatrix}
0\\
1
\end{pmatrix}
=
\begin{pmatrix}
0\cdot 0\\
0\cdot 1\\
1\cdot 0\\
1\cdot 1
\end{pmatrix}
=
\begin{pmatrix}
0\\
0\\
0\\
1
\end{pmatrix}
\]

\medskip
\noindent
したがって，
\[
|0\rangle\otimes|0\rangle=
\begin{pmatrix}
1\\0\\0\\0
\end{pmatrix},\quad
|0\rangle\otimes|1\rangle=
\begin{pmatrix}
0\\1\\0\\0
\end{pmatrix},\quad
|1\rangle\otimes|0\rangle=
\begin{pmatrix}
0\\0\\1\\0
\end{pmatrix},\quad
|1\rangle\otimes|1\rangle=
\begin{pmatrix}
0\\0\\0\\1
\end{pmatrix}.
\]
これらはちょうど $\mathbb{C}^4$ の標準基底
\[
e_1=
\begin{pmatrix}
1\\0\\0\\0
\end{pmatrix},\quad
e_2=
\begin{pmatrix}
0\\1\\0\\0
\end{pmatrix},\quad
e_3=
\begin{pmatrix}
0\\0\\1\\0
\end{pmatrix},\quad
e_4=
\begin{pmatrix}
0\\0\\0\\1
\end{pmatrix}
\]
そのものである。

\medskip
\noindent
さらに，任意のベクトル
\[
v=
\begin{pmatrix}
a\\
b\\
c\\
d
\end{pmatrix}
\in \mathbb{C}^4
\]
は
\[
v
=
a
\begin{pmatrix}
1\\0\\0\\0
\end{pmatrix}
+b
\begin{pmatrix}
0\\1\\0\\0
\end{pmatrix}
+c
\begin{pmatrix}
0\\0\\1\\0
\end{pmatrix}
+d
\begin{pmatrix}
0\\0\\0\\1
\end{pmatrix}
\]
と書けるので，
\[
v
=
a(|0\rangle\otimes|0\rangle)
+b(|0\rangle\otimes|1\rangle)
+c(|1\rangle\otimes|0\rangle)
+d(|1\rangle\otimes|1\rangle)
\]
が成り立つ．よって，これら 4 つのベクトルは $H$ を張る．

\medskip
\noindent
また，
\[
\alpha(|0\rangle\otimes|0\rangle)
+\beta(|0\rangle\otimes|1\rangle)
+\gamma(|1\rangle\otimes|0\rangle)
+\delta(|1\rangle\otimes|1\rangle)
=0
\]
と仮定すると，
\[
\alpha
\begin{pmatrix}
1\\0\\0\\0
\end{pmatrix}
+\beta
\begin{pmatrix}
0\\1\\0\\0
\end{pmatrix}
+\gamma
\begin{pmatrix}
0\\0\\1\\0
\end{pmatrix}
+\delta
\begin{pmatrix}
0\\0\\0\\1
\end{pmatrix}
=
\begin{pmatrix}
0\\0\\0\\0
\end{pmatrix}
\]
であるから、
\[
\alpha=\beta=\gamma=\delta=0
\]
を得る．したがって，これら 4 つのベクトルは一次独立である．

\medskip
\noindent
以上より，
\[
|0\rangle\otimes|0\rangle,\quad
|0\rangle\otimes|1\rangle,\quad
|1\rangle\otimes|0\rangle,\quad
|1\rangle\otimes|1\rangle
\]
は $H=\mathbb{C}^2\otimes\mathbb{C}^2\cong\mathbb{C}^4$ の標準基底である．
\end{proof}

\newpage

%============================================================
\section*{\boxed{\textbf{\large 問3}}}
%============================================================

Show that
\[
(I_2-\sigma_k)\otimes(I_2+\sigma_k)+(I_2+\sigma_k)\otimes(I_2-\sigma_k)=I_4-\sigma_k\otimes\sigma_k
\]
for $k=1,2,3$, where $I_n$ is the identity matrix in $M_n(\mathbb{C})$.

\vspace{1em}

%============================================================
\section*{\boxed{\textbf{\large 解答}}}
%============================================================

\begin{proof} 
\ \\
$k\in\{1,2,3\}$ を任意に固定する。
テンソル積は各変数について線形であるから，
\[
(A+B)\otimes C=A\otimes C+B\otimes C,\qquad
A\otimes(B+C)=A\otimes B+A\otimes C
\]
が成り立つ。また、
\[
(\lambda A)\otimes B=\lambda(A\otimes B)=A\otimes(\lambda B)
\qquad
(\lambda\in\mathbb{C})
\]
も成り立つ。これを用いて左辺を順に展開する。

\medskip
\noindent
まず、第 1 項は
\begin{align*}
(I_2-\sigma_k)\otimes(I_2+\sigma_k)
&=
I_2\otimes I_2+I_2\otimes\sigma_k-\sigma_k\otimes I_2-\sigma_k\otimes\sigma_k
\end{align*}
同様に、第 2 項は
\begin{align*}
(I_2+\sigma_k)\otimes(I_2-\sigma_k)
&=
I_2\otimes I_2-I_2\otimes\sigma_k+\sigma_k\otimes I_2-\sigma_k\otimes\sigma_k
\end{align*}

\medskip
\noindent
したがって、両者を足し合わせると
\begin{align*}
&(I_2-\sigma_k)\otimes(I_2+\sigma_k)+(I_2+\sigma_k)\otimes(I_2-\sigma_k)\\
&=
\bigl(I_2\otimes I_2+I_2\otimes\sigma_k-\sigma_k\otimes I_2-\sigma_k\otimes\sigma_k\bigr)\\
&\qquad
+\bigl(I_2\otimes I_2-I_2\otimes\sigma_k+\sigma_k\otimes I_2-\sigma_k\otimes\sigma_k\bigr)
\end{align*}

\medskip
\noindent
右辺で同種項をまとめると、
\begin{align*}
&(I_2-\sigma_k)\otimes(I_2+\sigma_k)+(I_2+\sigma_k)\otimes(I_2-\sigma_k)\\
&=
(I_2\otimes I_2)+(I_2\otimes I_2)
+(I_2\otimes\sigma_k)-(I_2\otimes\sigma_k)\\
&\qquad
-(\sigma_k\otimes I_2)+(\sigma_k\otimes I_2)
-(\sigma_k\otimes\sigma_k)-(\sigma_k\otimes\sigma_k)\\
&=
2(I_2\otimes I_2)-2(\sigma_k\otimes\sigma_k)
\end{align*}

\medskip
\noindent
ここで、単位行列のテンソル積について
\[
I_2\otimes I_2=
\begin{pmatrix}
I_2 & \mathbf{0}\\
\mathbf{0} & I_2
\end{pmatrix}
=
\begin{pmatrix}
1 & 0 & 0 & 0\\
0 & 1 & 0 & 0\\
0 & 0 & 1 & 0\\
0 & 0 & 0 & 1
\end{pmatrix}
=I_4
\]
であるから、
\[
2(I_2\otimes I_2)-2(\sigma_k\otimes\sigma_k)
=
2I_4-2\sigma_k\otimes\sigma_k
\]
を得る。よって
\[
(I_2-\sigma_{k})\otimes(I_2+\sigma_{k})+(I_2+\sigma_{k})\otimes (I_2-\sigma_{k})=2I_4-2\sigma_{k}\otimes\sigma_{k}
\]
が示された。

\medskip
\noindent
さらに、$k=1,2,3$ の各場合について確認すると、

\paragraph{\textbf{$k=1$ の場合}} \ \\
\[
\sigma_1=
\begin{pmatrix}
0 & 1\\
1 & 0
\end{pmatrix}
\]
であるから、
\[
I_2-\sigma_1=
\begin{pmatrix}
1 & -1\\
-1 & 1
\end{pmatrix},
\qquad
I_2+\sigma_1=
\begin{pmatrix}
1 & 1\\
1 & 1
\end{pmatrix}
\]
となる。したがって、
\[
(I_2-\sigma_1)\otimes(I_2+\sigma_1)
=
\begin{pmatrix}
1 & 1 & -1 & -1\\
1 & 1 & -1 & -1\\
-1 & -1 & 1 & 1\\
-1 & -1 & 1 & 1
\end{pmatrix}
\]
および
\[
(I_2+\sigma_1)\otimes(I_2-\sigma_1)
=
\begin{pmatrix}
1 & -1 & 1 & -1\\
-1 & 1 & -1 & 1\\
1 & -1 & 1 & -1\\
-1 & 1 & -1 & 1
\end{pmatrix}
\]
である。よって、
\[
(I_2-\sigma_1)\otimes(I_2+\sigma_1)
+(I_2+\sigma_1)\otimes(I_2-\sigma_1)
=
\begin{pmatrix}
2 & 0 & 0 & -2\\
0 & 2 & -2 & 0\\
0 & -2 & 2 & 0\\
-2 & 0 & 0 & 2
\end{pmatrix}.
\]

一方で、
\[
\sigma_1\otimes\sigma_1
=
\begin{pmatrix}
0 & 0 & 0 & 1\\
0 & 0 & 1 & 0\\
0 & 1 & 0 & 0\\
1 & 0 & 0 & 0
\end{pmatrix}
\]
であるから、
\[
2I_4-2\sigma_1\otimes\sigma_1
=
\begin{pmatrix}
2 & 0 & 0 & -2\\
0 & 2 & -2 & 0\\
0 & -2 & 2 & 0\\
-2 & 0 & 0 & 2
\end{pmatrix}.
\]
したがって、
\[
(I_2-\sigma_1)\otimes(I_2+\sigma_1)
+(I_2+\sigma_1)\otimes(I_2-\sigma_1)
=
2I_4-2\sigma_1\otimes\sigma_1
\]
が成り立つ。

\paragraph{\textbf{$k=2$ の場合}} \ \\
\[
\sigma_2=
\begin{pmatrix}
0 & -i\\
i & 0
\end{pmatrix}
\]
であるから、
\[
I_2-\sigma_2=
\begin{pmatrix}
1 & i\\
-i & 1
\end{pmatrix},
\qquad
I_2+\sigma_2=
\begin{pmatrix}
1 & -i\\
i & 1
\end{pmatrix}
\]
となる。したがって、
\[
(I_2-\sigma_2)\otimes(I_2+\sigma_2)
=
\begin{pmatrix}
1 & -i & i & 1\\
i & 1 & -1 & i\\
-i & -1 & 1 & -i\\
1 & -i & i & 1
\end{pmatrix}
\]
および
\[
(I_2+\sigma_2)\otimes(I_2-\sigma_2)
=
\begin{pmatrix}
1 & i & -i & 1\\
-i & 1 & -1 & -i\\
i & -1 & 1 & i\\
1 & i & -i & 1
\end{pmatrix}
\]
である。よって、
\[
(I_2-\sigma_2)\otimes(I_2+\sigma_2)
+(I_2+\sigma_2)\otimes(I_2-\sigma_2)
=
\begin{pmatrix}
2 & 0 & 0 & 2\\
0 & 2 & -2 & 0\\
0 & -2 & 2 & 0\\
2 & 0 & 0 & 2
\end{pmatrix}.
\]

一方で、
\[
\sigma_2\otimes\sigma_2
=
\begin{pmatrix}
0 & 0 & 0 & -1\\
0 & 0 & 1 & 0\\
0 & 1 & 0 & 0\\
-1 & 0 & 0 & 0
\end{pmatrix}
\]
であるから、
\[
2I_4-2\sigma_2\otimes\sigma_2
=
\begin{pmatrix}
2 & 0 & 0 & 2\\
0 & 2 & -2 & 0\\
0 & -2 & 2 & 0\\
2 & 0 & 0 & 2
\end{pmatrix}.
\]
したがって、
\[
(I_2-\sigma_2)\otimes(I_2+\sigma_2)
+(I_2+\sigma_2)\otimes(I_2-\sigma_2)
=
2I_4-2\sigma_2\otimes\sigma_2
\]
が成り立つ。

\paragraph{\textbf{$k=3$ の場合}} \ \\
\[
\sigma_3=
\begin{pmatrix}
1 & 0\\
0 & -1
\end{pmatrix}
\]
であるから、
\[
I_2-\sigma_3=
\begin{pmatrix}
0 & 0\\
0 & 2
\end{pmatrix},
\qquad
I_2+\sigma_3=
\begin{pmatrix}
2 & 0\\
0 & 0
\end{pmatrix}
\]
となる。したがって、
\[
(I_2-\sigma_3)\otimes(I_2+\sigma_3)
=
\begin{pmatrix}
0 & 0 & 0 & 0\\
0 & 0 & 0 & 0\\
0 & 0 & 4 & 0\\
0 & 0 & 0 & 0
\end{pmatrix}
\]
および
\[
(I_2+\sigma_3)\otimes(I_2-\sigma_3)
=
\begin{pmatrix}
0 & 0 & 0 & 0\\
0 & 4 & 0 & 0\\
0 & 0 & 0 & 0\\
0 & 0 & 0 & 0
\end{pmatrix}
\]
である。よって、
\[
(I_2-\sigma_3)\otimes(I_2+\sigma_3)
+(I_2+\sigma_3)\otimes(I_2-\sigma_3)
=
\begin{pmatrix}
0 & 0 & 0 & 0\\
0 & 4 & 0 & 0\\
0 & 0 & 4 & 0\\
0 & 0 & 0 & 0
\end{pmatrix}.
\]

一方で、
\[
\sigma_3\otimes\sigma_3
=
\begin{pmatrix}
1 & 0 & 0 & 0\\
0 & -1 & 0 & 0\\
0 & 0 & -1 & 0\\
0 & 0 & 0 & 1
\end{pmatrix}
\]
であるから、
\[
2I_4-2\sigma_3\otimes\sigma_3
=
\begin{pmatrix}
0 & 0 & 0 & 0\\
0 & 4 & 0 & 0\\
0 & 0 & 4 & 0\\
0 & 0 & 0 & 0
\end{pmatrix}.
\]
したがって、
\[
(I_2-\sigma_3)\otimes(I_2+\sigma_3)
+(I_2+\sigma_3)\otimes(I_2-\sigma_3)
=
2I_4-2\sigma_3\otimes\sigma_3
\]
が成り立つ。

\medskip
\noindent
以上より、$k=1,2,3$ のいずれの場合にも
\[
(I_2-\sigma_k)\otimes(I_2+\sigma_k)
+(I_2+\sigma_k)\otimes(I_2-\sigma_k)
=
2I_4-2\sigma_k\otimes\sigma_k
\]
が成り立つ。

\end{proof}


%============================================================
\appendix
\section*{\boxed{\textbf{\large Appendix }}}
%============================================================
\begin{prop}\label{proof:dim}
$\mathbb{F}$ を体とし、$V,W$ を $\mathbb{F}\in \{\mathbb{R}, \mathbb{C}\}$ 上の有限次元ベクトル空間とする。
このとき、次は同値である：
\[
V \cong W \qquad 
\qquad \Longleftrightarrow \qquad
\dim V = \dim W.
\]
\end{prop}

\begin{proof}
\ \\
$\mathbf{(1)}$\quad $(\Rightarrow)$ を示す。
$V \cong W$ であると仮定する。
すると、ある線形同型写像
\[
T:V\to W
\]
が存在する。

\medskip
\noindent
いま、$\dim V=n$ とおき、$V$ の基底を
\[
\{v_1,v_2,\dots,v_n\}
\]
とする。
このとき、$W$ におけるベクトルの族
\[
\{T(v_1),T(v_2),\dots,T(v_n)\}
\]
を考える。
これが $W$ の基底であることを示せば、
\[
\dim W=n=\dim V
\]
が従う。

\paragraph{$\{T(v_1),T(v_2),\dots,T(v_n)\}$ \textbf{が} $W$ \textbf{を張ることを示す。}} \ \\
任意に $w\in W$ をとる。
$T$ は全射であるから、ある $v\in V$ が存在して
\[
T(v)=w
\]
となる。\\
\noindent
ここで、$\{v_1,v_2,\dots,v_n\}$ は $V$ の基底であるから、ある係数
\[
a_1,a_2,\dots,a_n\in\mathbb{F}
\]
を用いて
\[
v=a_1v_1+a_2v_2+\cdots+a_nv_n
\]
と一意的に表せる。
よって、$T$ の線形性より
\[
w=T(v)
=T(a_1v_1+a_2v_2+\cdots+a_nv_n)
=a_1T(v_1)+a_2T(v_2)+\cdots+a_nT(v_n)
\]
となる。
したがって、任意の $w\in W$ は $\{T(v_1),T(v_2),\dots,T(v_n)\}$ の線形結合として表される。
ゆえに、
\[
W=\mathrm{span}\{T(v_1),T(v_2),\dots,T(v_n)\}
\]
である。

\paragraph{$\{T(v_1),T(v_2),\dots,T(v_n)\}$ \textbf{が一次独立であることを示す。}} \ \\
\[
b_1T(v_1)+b_2T(v_2)+\cdots+b_nT(v_n)=0
\qquad (b_1,\dots,b_n\in\mathbb{F})
\]
と仮定する。
このとき、$T$ の線形性により
\[
T(b_1v_1+b_2v_2+\cdots+b_nv_n)=0
\]
を得る。\\
\noindent
$T$ は単射であるから、
\[
b_1v_1+b_2v_2+\cdots+b_nv_n=0
\]
である。
ところが、$\{v_1,v_2,\dots,v_n\}$ は基底、一次独立であるから、
\[
b_1=b_2=\cdots=b_n=0
\]
が従う。
したがって、$\{T(v_1),T(v_2),\dots,T(v_n)\}$ は一次独立である。

\medskip
\noindent
以上より、$\{T(v_1),T(v_2),\dots,T(v_n)\}$ は $W$ の基底である。
したがって、
\[
\dim W=n=\dim V
\]
となる。
よって、
\[
V\cong W \Longrightarrow \dim V=\dim W
\]
が示された。

\vspace{1em}

\noindent
$\mathbf{(2)}$\quad $(\Leftarrow)$ を示す。
\[
\dim V=\dim W
\]
と仮定する。
この共通の次元を $n$ と書くと、
\[
\dim V=\dim W=n.
\]

\medskip
\noindent
$V$ と $W$ は有限次元であるから、それぞれ基底をとることができる。\\
\noindent
$V$ の基底を
\[
\{v_1,v_2,\dots,v_n\},
\]
$W$ の基底を
\[
\{w_1,w_2,\dots,w_n\}
\]
とする。\\
\noindent
このとき、写像 $T:V\to W$ を
\[
T(v_i)=w_i\qquad (i=1,2,\dots,n)
\]
で定め、任意の$v\in V$に対して
\[
v=a_1v_1+a_2v_2+\cdots+a_nv_n
\]
と書ける。このとき、
\[
T(v):=a_1w_1+a_2w_2+\cdots+a_nw_n
\]
となる。

% \paragraph{\textbf{この定義が well-defined であることの確認}} \ \\
% $\{v_1,v_2,\dots,v_n\}$ は $V$ の基底であるから、任意の $v\in V$ の表示
% \[
% v=a_1v_1+a_2v_2+\cdots+a_nv_n
% \]
% は一意的である。
% したがって、上の式による $T(v)$ の定義は表し方に依らず一意に定まる。
% よって、$T$ は well-defined である。

\paragraph{\textbf{$T$ が線形であることの確認}} \ \\
任意の
\[
v=a_1v_1+\cdots+a_nv_n,\qquad
u=b_1v_1+\cdots+b_nv_n
\]
および任意の $\alpha,\beta\in\mathbb{F}$ に対して、
\[
\alpha v+\beta u
=
(\alpha a_1+\beta b_1)v_1+\cdots+(\alpha a_n+\beta b_n)v_n
\]
であるから、
\begin{align*}
T(\alpha v+\beta u)
&=
(\alpha a_1+\beta b_1)w_1+\cdots+(\alpha a_n+\beta b_n)w_n\\
&=
\alpha(a_1w_1+\cdots+a_nw_n)+\beta(b_1w_1+\cdots+b_nw_n)\\
&=
\alpha T(v)+\beta T(u).
\end{align*}
したがって、$T$ は線形写像である。

\paragraph{\textbf{$T$ が単射であることを示す。}} \ \\
$T(v)=0$ と仮定する。
$v\in V$ を
\[
v=a_1v_1+a_2v_2+\cdots+a_nv_n
\]
と書けば、
\[
0=T(v)=a_1w_1+a_2w_2+\cdots+a_nw_n
\]
となる。
ところが、$\{w_1,w_2,\dots,w_n\}$ は $W$ の基底であるから一次独立である。
したがって、
\[
a_1=a_2=\cdots=a_n=0
\]
である。
よって、
\[
v=0
\]
となる。
したがって、$\ker T=\{0\}$ である。\\
\noindent
さらに、$u,v\in V$ に対して、
$T(u) = T(v)$ とすると、$T$の線形性より、$T(u-v) = 0$ となる。
したがって、$u-v \in \ker T$ であり、$\ker T=\{0\}$ であるから、$u-v=0$ すなわち $u=v$ となる。ゆえに、$T$ は単射である。

\paragraph{\textbf{$T$ が全射であることを示す。}} \ \\
任意に $w\in W$ をとる。
$\{w_1,w_2,\dots,w_n\}$ は $W$ の基底であるから、ある係数
\[
c_1,c_2,\dots,c_n\in\mathbb{F}
\]
が存在して
\[
w=c_1w_1+c_2w_2+\cdots+c_nw_n
\]
と書ける。
ここで、
\[
v=c_1v_1+c_2v_2+\cdots+c_nv_n \in V
\]
とおけば、$T$ の定義より
\[
T(v)=c_1w_1+c_2w_2+\cdots+c_nw_n=w  
\]
となる。
したがって、任意の $w\in W$ は $T$ の像に属する。\\
\noindent
\ie $\forall w \in W, w\in \mathrm{Im}(T) = \{T(v) |v \in V\}$が成り立つから、$W \subset \mathrm{Im}(T)$である。定義により、$\mathrm{Im}(T) \subset W$ であるから、$W = \mathrm{Im}(T)$ となる。ゆえに、$T$ は全射である。\\[0.5em]
以上より、$T$ は線形かつ単射かつ全射である。
したがって、$T$ は線形同型写像である。
ゆえに、
\[
V\cong W
\]
が従う。
すなわち、
\[
\dim V=\dim W \Longrightarrow V\cong W
\]
である。\\[1.5em]
以上の$\mathbf{(1)}$と$\mathbf{(2)}$で、
\[
V \cong W
\qquad \Longleftrightarrow \qquad
\dim V=\dim W
\]
が示された。
\end{proof}

\end{document}